\scanpage{53}

\deleted{(i)} et \unsure{considérons}

(i) \diagram{\begin{tikzcd}
W_{\mathcal M}^{-1}\mathcal M\arrow[r,"{\deleted{\unsure{\varphi}}\overline\lambda_!}"]&\mathrm D_\bullet(A)
\end{tikzcd}}

\margin{Indépendamment des \unsure{a) et b) b')}
est-il vrai que si $\overline{\lambda^*}$ existe, alors
$\overline\lambda_!$ et $\overline\lambda_* $ adjoints l'un de l'autre ??}

Conditions équivalentes

\unsure{(i)} \deleted{$\lambda_!$ et $\lambda^*$ \unsure{vérifient}}
$\lambda^{*-1}(W_{\mathcal M}^{-1})=W_{\mathrm{Ch}_\bullet(A)}$,
et les foncteurs déduits par passage aux cat. localisées

\diagram{\begin{tikzcd}[column sep=large]
W_{\mathcal M}^{-1}\mathcal M\arrow[r,shift left,"\overline\lambda_!"]&
\mathrm D_\bullet(A)=W_{\mathrm{Ch}_\bullet(A)}^{-1}\mathrm{Ch}_\bullet(A)\arrow[l,shift left,"\overline{\lambda^*}"]
\end{tikzcd}}

sont quasi-inv. l'un de l'autre, les isom. d'adjonction déduits de ceux pour $\lambda_!$, $\lambda^*$

\unsure{(ii)} $\forall\ l_\bullet\in\operatorname{Ob}\mathrm D_\bullet(A)$, ---
\[
\lambda_!\lambda^*(l_\bullet)\longrightarrow l_\bullet\quad\text{qu. iso.}
\]

\unsure{(ii')} Comme \unsure{(ii)}, mais avec $l_\bullet$ supposé \uline{projectif}
i.e.
\[
\left\{\begin{array}{ll}
\unsure{ii'_a)}&l_i\text{ \inserted{projectif}}\ \deleted{\unsure{\text{complexe}}}\quad\forall i\in\mathbb N\\
\unsure{ii'_b)}&H_i(l_\bullet)=0\quad\text{si }i>0\\
\unsure{ii'_b)|ii'_c)}&H_0(l_\bullet)\text{ projectif}
\end{array}\right.
\]

\unsure{(ii'')} Comme \unsure{(ii')} mais avec $l_\bullet$ dans un des deux cas

($\alpha$) $l_\bullet$ réduit au degré zéro, projectif

($\beta$) $l_\bullet$ à composantes projectives,
\deleted{\unsure{$H_k$}} et qu. iso à $0$

\uline{Lemme 3} Supposons $\mathcal M$ \inserted{\uline{$k$-abélienne}}
munie d'une structure multiplicative $\otimes$ \inserted{\unsure{$k$-bilinéaire}}
ass. comm. unit., \deleted{Support} unité $1_{\mathcal M}$
\inserted{et $\varphi$ adjoint à gauche de
$\psi:A\longrightarrow\mathcal M$, $H\longmapsto H\otimes_k1_{\mathcal M}$}
\margin{$A=Ab_k$}

Supposons
\[
\left\{\begin{array}{ll}
\unsure{(i)}&L_i\varphi(1_{\mathcal M})=0\quad\text{pour }i>0\\[4pt]
\unsure{(ii)}&\exists\quad 1_{\mathcal M}\lhook\joinrel\xrightarrow{\alpha}J,
\quad\text{tel que}
\end{array}\right.
\]
\margin{i.e. $\lambda_!(1_{\mathcal M})$ est une \uline{résolution} de $\varphi(1_{\mathcal M})$}

\[
\left\{\begin{array}{ll}
\unsure{ii'_1)}&\alpha\otimes\mathrm{id}_x\text{ mono}\quad\forall x\text{ dans }\mathcal M\\
\unsure{ii'_2)}&L\varphi(J\otimes x)=0
\end{array}\right.
\]
\margin{\deleted{\unsure{et}} (plus fort que $L_i\varphi(J\otimes x)=0$ pour $i>0$ :
on le veut pour $i=0$ aussi)}

Alors \deleted{\unsure{toutes}} les conditions du lemme 2 sont satisfaites

\mspage{52}
